\section{Related work}\label{sec:related-work}

\subsection{Physics of vibration-driven
locomotion}

A bristlebot converts vertical vibration into horizontal motion through
the asymmetric friction of tilted legs. Cicconofri and DeSimone
\cite{cicconofri2015motility} derive the mean speed as a function of
vibration frequency for a model bristle-bot in closed form, including
the frequency at which the direction of motion reverses, and the
reversal was demonstrated experimentally on a vibrated substrate
\cite{cicconofri2016inversion}. Kim et al. \cite{kim2020millibristle}
model milli-bristle-bots with dry friction and predict separate forward
and backward frequency bands and resonances, confirmed in experiment.
Giomi et al. \cite{giomi2013bristlebots} show that a single per-robot
number, the ratio of spinning to walking, selects the collective state
of confined bristlebots, which is the clearest argument that per-robot
calibration matters at the collective level. Scholz et al.
\cite{scholz2018inertial} measure a lag between orientation and
velocity in vibrobots and explain it with an underdamped active-Brownian
model, which motivates fitting a first-order lag rather than assuming an
instantaneous response. For two-motor platforms, Vartholomeos and
Papadopoulos \cite{vartholomeos2006microrobot} analysed a planar
micro-robot driven by two centripetal-force actuators, and Notomista et
al. \cite{notomista2019brushbot} derived brushbot dynamics from first
principles, with differential-drive kinematics \(v=(v_L+v_R)/2\),
\(\omega=(v_R-v_L)/W\) in which each \(v_i\) is nonlinear in motor
speed; their experiments validate the fully actuated variant only. Two
consequences for us: a monotone per-motor activation can fail at high
drive where the reversal or a resonance sits, so the top of the drive
range must be checked before the model form is fixed; and the
differential structure is a hypothesis to test.

\subsection{Characterising and calibrating vibration
robots}

The closest precedent for an automated calibration loop is the Droplet
platform \cite{klingner2014droplet}, whose stick-slip vibration-motor
powertrain is calibrated per robot by a Nelder--Mead search over motor
parameters, with commands sent over infrared and paths tracked by a
webcam. It searches per robot and, as far as its abstract states, fits
no model. Swarmodroid
\cite{dmitriev2023swarmodroid} characterises single-motor bristle-bots
with ArUco tracking: speed rises monotonically with PWM, the vibration
frequency rises roughly linearly with PWM, and the dispersion of speed
grows with PWM, for seven robots of each type; there is no two-motor
map, inverse or population model. Pogobot \cite{loi2025pogobot}, the
platform closest to GRASPion in purpose, is characterised with an
overhead camera on five robots and ten trials per PWM level with equal
power on both motors, for three locomotion modules: wheels (up to about
25 cm/s, the ``most consistent dynamics across individuals''),
uncalibrated vibration legs (up to about 7 cm/s with a stronger turning
bias) and brushes that mostly spin; its motor calibration is an onboard
yaw-rate loop that trims the command asymmetry until the robot drives
straight, applied to the wheeled units only, with the vibration units
explicitly left uncalibrated; we found no variance numbers in it.
Patterson and Parisi \cite{patterson2023fundamental} take the distinct
``free speeds'' of vibration-driven vehicles as a given input to a
traffic study rather than something to calibrate. For GRASPion's
predecessor, Noirhomme et al. \cite{noirhomme2025brainbots} give the
most detailed motion characterisation (centre and long-axis tracking,
Savitzky--Golay differentiation, instantaneous centre of rotation, a
spin-versus-walk index swept over motor voltage and leg angle), and for
GRASPion itself the platform paper reports a single-motor rotation and a
diffusive mean-squared displacement with no drive dependence
\cite{novkoski2026graspion}. The digital-twin paper
\cite{mammadli2026twin} adds a camera-based \((v,\omega)\to(L,R)\)
calibration at a high level; its public code shows a 7×7 grid of powers
60--120 on one robot, one value per cell, a cubic / quadratic polynomial
fit inverted by bounded least squares into two 68×490 lookup tables, and
separate clockwise and counter-clockwise centre-of-rotation offsets
(14.5 and 16.1 mm). Our reanalysis of those tables (§\ref{sec:lab-tables}) finds the
cubic to overfit and only about 44 \% of the lookup table's box to be
reachable.

The hands-on Kilobot procedure sets left, right and straight motor values per robot and
stores them in EEPROM \cite{rubenstein2014kilobot}; the lessons-learned
section calls it tedious and hard to automate. ARK \cite{reina2017ark}
provides overhead tracking and infrared control for a hundred robots and
says automatic motor calibration is enabled but does not describe it.
Pinciroli et al. \cite{pinciroli2018kilobots} measured 120
pre-calibrated Kilobots in one-minute straight runs, estimated a
per-robot Gaussian speed bias (SD 1.86 mm/s) and injected it into the
ARGoS simulator; the population variance was measured, but used only as
simulator noise. Raoufi et al. \cite{raoufi2023individuality} show a
persistent per-robot heading bias whose within-robot spread is smaller
than the ensemble's, and that calibration is ``only temporary''. the Kilombo
simulator \cite{jansson2015kilombo} motivates itself partly by the need
for frequent recalibration; the Robotarium automates registration and sensor calibration of a remote
wheeled swarm with an overhead tracker \cite{pickem2016robotarium}.
Siemensma and Haghighat \cite{siemensma2024bayesian} model per-robot
misalignment and speed factors of five wheeled robots as hand-tuned
distributions for a simulator. In experimental active matter, per-unit
identification is routine: Siebers et al. \cite{siebers2023disorder}
fit twenty light-driven circle walkers individually as active Brownian
particles and exploit their disorder, and find a hard light-intensity
threshold below which the walkers do not move, a dead zone in our
language; Kiechl et al. \cite{kiechl2026chiral} show that a chiral
Hexbug matches the active-Brownian circle-swimmer predictions for the
mean-squared displacement and intermediate scattering function; Bailey
et al. \cite{bailey2022logit} show that ordinary least squares on the
mean-squared displacement is biased and recommend a log-MSD fit with
bootstrap intervals. What is missing across all of this is a measured map from
the two motor powers and the battery voltage to the body twist for any
two-motor vibration robot, a quantitative robot / day / segment
decomposition of its variability, and any use of a population of
calibrated robots to calibrate the next one.

\subsection{Kinematic calibration and identifiability of mobile
robots}

Calibrating a motion model from external pose measurements is a mature
problem for wheeled robots. Antonelli et al.
\cite{antonelli2005odometry} fit differential-drive parameters by
linear least squares from camera measurements on arbitrary paths; Censi
et al. \cite{censi2013calibration} give a closed-form
maximum-likelihood calibration of odometry and sensor parameters with an
observability analysis and show accuracy near the Cramér--Rao bound;
Kelly \cite{kelly2004odometry} calibrates systematic and stochastic
(diffusion) odometry parameters jointly from path integrals; and
Seegmiller et al. \cite{seegmiller2013ipem} fit the integrated model to
low-rate pose data by minimising multi-step prediction error, including
a first-order powertrain lag and delay, which is both an alternative to
fitting lag from step responses and the natural validation metric for a
motion model. Slip-dominated vehicles are modelled by instantaneous
centres of rotation that stay in a bounded region: Martínez et al.
\cite{martinez2005icr} fit constant ICR positions for tracked robots,
Mandow et al. \cite{mandow2007skidsteer} extend this to wheeled
skid-steer platforms on several surfaces, and Yi et al.
\cite{yi2009skidsteer} estimate ICR locations online with an extended
Kalman filter; the ICR positions are fitted as constants from
steady-state runs, and the degeneracy we meet in
§\ref{sec:identifiability} is of a different kind, between two offsets
that enter only through the sign of the yaw rate. Online self-calibration tracks parameters that drift: Roy and
Thrun \cite{roy1999selfcalibration} with incremental maximum likelihood
and Martinelli et al. \cite{martinelli2007selfcalibration} with an
augmented Kalman filter for the systematic part and a second filter for
the noise part. Becker et al. \cite{becker2014massive} steer many
differently calibrated robots with one broadcast command plus online
gain estimation. Our forward map is in the Antonelli--Martínez family,
with motor commands in place of wheel encoders and a nonlinear per-motor
activation in front of the linear twist. That activation is a
Hammerstein structure, a static nonlinearity followed by a linear block,
for which dead-zone models of DC motors are standard
\cite{kara2004dcmotor, zhang2017deadzone} and for which Bayesian
identification by empirical Bayes exists in nonparametric form
\cite{risuleo2017hammerstein}; we found no published fit of a dead zone
plus power law to a vibration motor's speed against PWM, so the form
itself is a hypothesis that the grid tests.

\subsection{Choosing the calibration
experiments}

Bayesian experimental design chooses measurements to maximise expected
information \cite{chaloner1995design, rainforth2024design}. For robot
calibration, Sun and Hollerbach \cite{sun2008observability} relate the
observability indices used in kinematic calibration to the alphabetic
optimality criteria and show that D-optimality minimises parameter
variance while E-type criteria minimise end-effector position
uncertainty. Active
learning of dynamics models chooses where to sample under the
constraints the dynamics impose \cite{buissonfenet2020active}, and
receding-horizon curiosity \cite{schultheis2019curiosity} frames
exploration as sequential Bayesian design alternating with nonlinear
system identification. Preiss et al. \cite{preiss2017selfcalibration}
add an observability cost to trajectory optimisation for visual-inertial
self-calibration, and Das et al. \cite{das2024boed} apply Bayesian
optimisation to the choice of calibration poses for a manipulator. For
wheeled odometry the choice of calibration path has been studied without
a Bayesian framing, for instance the size of the UMBmark square under
non-systematic errors \cite{jung2011tracks}, and fully automatic
wheel-and-camera calibration of Duckiebots from fiducials reaches the
accuracy of manual calibration \cite{chaika2021duckiebot}. Beyond such track-design studies we found, in a
non-systematic search, no Bayesian or alphabetic-optimal design of a
mobile robot's motor-map or odometry calibration, and no published curve
of calibration error against the number of calibration samples for
cheap or swarm robots on hardware.

\subsection{Learned dynamics and few-shot
adaptation}

Nagabandi et al. \cite{nagabandi2018millirobot} learn a neural forward
model of a legged millirobot without odometry from under seventeen
minutes of random commands tracked externally and use it for
model-predictive control; it is the closest published experiment to ours
in platform and data budget and the natural neural comparison for our
grey-box map. Meta-learning adapts a dynamics model from a few recent
steps: across terrains and damage on a millirobot
\cite{nagabandi2019adapt}, across suspended payloads on a quadrotor
\cite{belkhale2021payload}, and across randomised simulated parameters
for legged robots that then adapt on hardware in a fraction of a second
without real calibration data \cite{kumar2021rma}. Neural-Fly
\cite{oconnell2022neuralfly} learns a shared basis from twelve minutes
of flight on one drone and adapts only the linear coefficients that mix
it, and the basis transferred to a different drone with a modest loss;
ALPaCA \cite{harrison2018alpaca} meta-learns features and a Gaussian
prior on last-layer weights with closed-form online Bayesian updates;
PACOH \cite{rothfuss2021pacoh} shows that learned priors overfit when
few meta-training tasks are available, which is the regime of a small
robot fleet. Meta-learned state-space models
\cite{chakrabarty2023meta}, in-context identification over a class of
systems \cite{forgione2023incontext} and universal car dynamics models
\cite{xiao2024anycar} all learn the across-unit distribution from many
systems, simulated or real. Xiao et al. \cite{xiao2021ikd} learn the
inverse map directly from driving data, and RAPTOR
\cite{eschmann2025raptor} distils a single recurrent policy from a
thousand simulated quadrotors that adapts in context on ten real ones,
handling fleet heterogeneity implicitly rather than through per-robot
parameters. Our model has the Neural-Fly / ALPaCA structure, a physical
basis with per-unit coefficients, except that the two dead-zone and
exponent parameters per motor are per-unit and nonlinear too; with
thirteen interpretable parameters a hierarchical Gaussian prior replaces
the neural network.

\subsection{Transfer across nominally identical
units}\label{sec:rw-transfer}

Population-based structural health monitoring treats a fleet of
nominally identical structures as a population whose shared behaviour is
a ``form'' and whose between-member variation is the population variance
\cite{bull2021pbshm}; hierarchical Bayesian multitask models over a
truck fleet and a wind farm let data-poor groups borrow strength from
data-rich ones and tighten interpretable parameters
\cite{bull2022fleet}, and the same machinery applies to a small
population of nominally identical blades \cite{dardeno2024frf}. In
system identification, hierarchical Bayesian models draw each dataset's
ODE parameters from a shared population distribution and, on simulated
sparse data, halve the trajectory error of unpooled fits
\cite{brugnara2026hierarchical}; PAC-Bayesian meta-learning learns a
prior over linear system matrices from training systems and studies its
sensitivity to their number \cite{huang2026pacbayes}. The closest
quantitative answer to ``how many units does a population prior need''
comes from battery research: Dechent et al.
\cite{dechent2021samplesize} estimate that one-, two- and
three-parameter ageing models need at least 9, 11 and 13 cells to
resolve cell-to-cell variability, and Zhou and Howey
\cite{zhou2023battery} separate cell-level manufacturing variability
from population-level cycling conditions in a hierarchical lifetime
model. In robotics, Williams et al. \cite{williams2008mtgp} learn
inverse dynamics as a multi-task Gaussian process whose task similarity
comes from inertial parameters and beat both independent and pooled
models; Bocsi et al. \cite{bocsi2013alignment} align data across robot
architectures. Helwa and Schoellig \cite{helwa2017transfer} transfer
six minutes of source-drone data to a target drone through a dynamic
(not static) map identified from simple experiments, and Sorocky et al.
\cite{sorocky2020selection} show with several candidate source robots
that the choice of source changes target error by 62 \% and that a
poorly matched source gives negative transfer, the clearest evidence we
found against relying on a single reference robot. In analytical chemistry the
same problem is called calibration transfer or instrument
standardisation
\cite{wang1991standardization, feudale2002transfer, workman2018transfer}:
a master instrument's calibration is transferred to a slave through a
small set of transfer samples measured on both, and this master/slave
form is the baseline we call \texttt{ref\_gain} in
§\ref{sec:baselines}.
